\section[Five minimal representations of the $\cR$-function]{Five minimal representations of the $\boldsymbol{\cR}$-function}\label{section3}

We begin this section by focusing on the kernel function
\begin{gather}
K(b;x,v):=G((\pm x \pm v -ib)/2)\equiv\prod_{\de_1,\de_2=+,-}G((\de_1 x+\de_2 v-ib)/2).
\end{gather}
We have established that this function satisf\/ies three independent kernel identities. We expect that these might be useful in other contexts than the present one. Indeed, here we only need the special case~\eqref{dv} of the f\/irst of the identities. We  collect the three identities in the following proposition.

\begin{prop}\label{proposition3.1}
Letting $b,d\in\C$, we have
\begin{gather}
\frac{s_{\de}(x-ib+id)}{s_{\de}(x)}K(b;x-ia_{-\de},v)+(i\to -i)\nonumber\\
\qquad{} =\frac{s_{\de}(v-id)}{s_{\de}(v)}K(b;x,v-ia_{-\de})+(i\to -i),\label{idd}
\\
\frac{s_{\de}(x-ib)}{s_{\de}(x)}K(b;x-2ia_{-\de},v)+(i\to -i)\nonumber\\
\qquad{}=
\frac{s_{\de}(v-ib)}{s_{\de}(v)}K(b;x,v-2ia_{-\de})+(i\to -i),\label{id2}
\\
\frac{s_{\de}((x-ib)/2)}{s_{\de}(x/2)}K(b;x-ia_{-\de},v)+(i\to -i)\nonumber\\
\qquad{} =
\frac{s_{\de}((v-ib)/2)}{s_{\de}(v/2)}K(b;x,v-ia_{-\de})+(i\to -i).\label{id3}
\end{gather}
\end{prop}

\begin{proof}
To prove \eqref{idd}, we divide l.h.s.\ and r.h.s.\ by $K(b;x-ia_{-\de},v)$ and use the A$\De$Es~\eqref{Gades} to write the result as
\begin{gather}
  \frac{s_{\de}(x-ib+id)}{s_{\de}(x)}+\frac{s_{\de}(x+ib-id)}{s_{\de}(x)}\frac{c_{\de}((x+ v-ib)/2)}{c_{\de}((x+ v+ib)/2)}\frac{c_{\de}((x- v-ib)/2)}{c_{\de}((x- v+ib)/2)}  \nonumber \\
\qquad{} = \frac{s_{\de}(v-id)}{s_{\de}(v)} \frac{c_{\de}((x-v-ib)/2)}{c_{\de}((x-v+ib)/2)} +   \frac{s_{\de}(v+id)}{s_{\de}(v)} \frac{c_{\de}((x+v-ib)/2)}{c_{\de}((x+v+ib)/2)}.\label{idq}
\end{gather}
Both sides are $2ia_{\de}$-periodic functions of $x$ with equal limits
\begin{gather}
e_{\de}(\pm (-ib+id))+e_{\de}(\pm (-ib-id)),\qquad \re x \to\pm\infty.
\end{gather}
The residues at the (generically simple) poles $x=0$, $x=ia_{\de}$ in the period strip clearly cancel. By Liouville's theorem, it remains to check that the residues at the poles $x=\pm v -ib/2\pm ia_{\de}$ cancel as well, and this is a routine calculation.

Next, we divide~\eqref{id2} by $K(b;x,v)$ and use~\eqref{Gades} to obtain
\begin{gather}
  \frac{s_{\de}(x-ib)}{s_{\de}(x)}\frac{c_{\de}((x-ia_{-\de}/2\pm v +ib)/2)}{c_{\de}((x-ia_{-\de}/2\pm v -ib)/2)}+ (x\to -x)
\nonumber \\
\qquad{} = \frac{s_{\de}(v-ib)}{s_{\de}(v)} \frac{c_{\de}((v-ia_{-\de}/2\pm x +ib)/2)}{c_{\de}((v-ia_{-\de}/2\pm x -ib)/2)}+ (v\to -v).
\end{gather}
Both sides are $2ia_{\de}$-periodic functions of $x$ with equal limits
\begin{gather}
e_{\de}(\pm ib)+e_{\de}(\mp ib)=\frac{s_{\de}(v-ib)}{s_{\de}(v)} +\frac{s_{\de}(v+ib)}{s_{\de}(v)},\qquad \re x \to\pm\infty.
\end{gather}
The residues at $x=0$ and $x=ia_{\de}$ manifestly cancel. It is a straightforward calculation to verify that the residues at the remaining poles $x=\pm ia_{\de}\pm ia_{-\de}/2\pm v +ib$ cancel, too. Hence~\eqref{id2} follows.

Finally, to prove~\eqref{id3}, we divide both sides by
$K(b;x-ia_{-\de},v)$ and use~\eqref{Gades} to get as the counterpart of~\eqref{idq}:
\begin{gather}
   \frac{s_{\de}((x-ib)/2)}{s_{\de}(x/2)}+\frac{s_{\de}((x+ib)/2)}{s_{\de}(x/2)}\frac{c_{\de}((x+ v-ib)/2)}{c_{\de}((x+ v+ib)/2)}\frac{c_{\de}((x- v-ib)/2)}{c_{\de}((x- v+ib)/2)}  \nonumber \\
\qquad{} = \frac{s_{\de}((v-ib)/2)}{s_{\de}(v/2)} \frac{c_{\de}((x-v-ib)/2)}{c_{\de}((x-v+ib)/2)} +   \frac{s_{\de}((v+ib)/2)}{s_{\de}(v/2)} \frac{c_{\de}((x+v-ib)/2)}{c_{\de}((x+v+ib)/2)}.
\end{gather}
Both sides are $2ia_{\de}$-periodic functions of $x$ with equal limits
\begin{gather}
e_{\de}(\pm (-ib/2))+e_{\de}(\pm (-3ib/2)),\qquad \re x \to\pm\infty.
\end{gather}
As before, residue cancellation at $x=0$ and $x=ia_{\de}$ is immediate, whereas the verif\/ication that the residues at the remaining poles $x=\pm ia_{\de} \pm v  -ib/2$ cancel as well involves a bit more work.
\end{proof}


From~\eqref{Apm} we see that the identity~\eqref{idd} can be rewritten as
\begin{gather}
A_{\de}(b-d;x)K(b;x,v)=A_{\de}(d;v)K(b;x,v),\qquad \de=+,-.
\end{gather}
At f\/irst sight, one might think that the identities~\eqref{id2} and~\eqref{id3} can also be rewritten by using a rescaled version of the two commuting $A_1$ dif\/ference operators $A_{\pm}(b;z)$. The two dif\/ference operators
\begin{gather}
\frac{s_{\de}((z-ib)/2)}{s_{\de}(z/2)}T^z_{ia_{-\de}} +
\frac{s_{\de}((z+ib)/2)}{s_{\de}(z/2)}T^z_{-ia_{-\de}},\qquad \de=+,-,
\end{gather}
featuring in~\eqref{id3}, do not even commute, however. Thus no such rescaling is possible for~\eqref{id3}. The dif\/ference operators
\begin{gather}
\frac{s_{\de}(z-ib)}{s_{\de}(z)}T^z_{2ia_{-\de}} +
\frac{s_{\de}(z+ib)}{s_{\de}(z)}T^z_{-2ia_{-\de}},\qquad \de=+,-,
\end{gather}
corresponding to~\eqref{id2} do commute. Even so, one can only rescale one of the operators such that it takes the~$A_1$ form~\eqref{Apm}, but not both at once.

For the purpose of studying the $A_1$ operators, then, we can only make use of~\eqref{idd}. More specif\/ically, our starting point is the special case $d=0$:
\begin{gather}\label{dv}
A_{\de}(b;x)K(b;x,v)= \big(T^v_{ia_{-\de}} +T^v_{-ia_{-\de}} \big)K(b;x,v).
\end{gather}
In order to exploit this identity, we introduce  the Fourier transform
\begin{gather}\label{defB}
B(b;x,y)\equiv \frac{1}{2}\int_{\R}dvK(b;x,v)\exp(i\alpha vy/2),\qquad b\in(0,2a),\qquad x,y\in \R.
\end{gather}
The integral is well def\/ined, since the $b$-restriction ensures that the $v$-poles of the integrand at
\begin{gather}\label{Kpo}
\pm v=x+2ia-ib+z_{kl}, \qquad \pm v=-x+2ia-ib+z_{kl},\qquad k,l\in\N,
\end{gather}
(cf.~\eqref{Gpo}--\eqref{zkl}), stay away from the real axis, and since the $G$-asymptotics~\eqref{Gas} entails an exponential decay
\begin{gather}
K(b;x,v)\sim \exp(\mp \alpha bv/2),\qquad b>0,\qquad x\in\C,\qquad \re v\to\pm\infty.
\end{gather}
These features also imply that $B$ extends from the real $x$-axis to a function that is holomorphic in the strip $\im x\in (-2a+b,2a-b)$.

Next, we temporarily  assume
\begin{gather}
b\in(0,a_s/2),\qquad a_s\equiv\min(\aaa).
\end{gather}
Then the action of the shifts in the A$\De$Os $A_{\pm}(b;x)$ given by~\eqref{Apm} is well def\/ined on $B(b;x,y)$, provided we restrict~$x$ to a strip $|\im x|<a_s/2$. Moreover, we may take the shifts under the integral sign in~\eqref{defB} and use the kernel identity~\eqref{dv} to obtain
\begin{gather}
A_{\de}(b;x)B(b;x,y)=\frac{1}{2}\int_{\R}dv\sum_{\tau=+,-}K(b;x,v+\tau ia_{-\de})\exp(i\pi vy/a_+a_-),\\
 |\im x|<a_s/2.\nonumber
\end{gather}
Upon shifting contours $\R\to \R\pm ia_{-\de}$, no poles are met, and so we obtain
\begin{gather}
\frac12\sum_{\tau=+,-}e_{\de}(\tau y)\int_{\R-i\tau a_{-\de}}dvK(b;x,v)\exp(i\pi vy/a_+a_-).
\end{gather}
The integrands of both terms are now equal, and the contours can be shifted back to $\R$ without changing the value of the integrals. Hence we deduce the eigenvalue equations
\begin{gather}\label{eigen}
A_{\pm}(b;x)B(b;x,y)=2c_{\pm}(y)B(b;x,y),\qquad |\im x|<a_s/2.
\end{gather}

Reverting to our previous assumption $b\in(0,2a)$, we proceed to obtain two dif\/ferent representations of $B(b;x,y)$. To this end we use the Plancherel relation
\begin{gather}\label{Plan}
\int_{\R}dpf(p)g(p)=\frac{\alpha}{2\pi}\int_{\R}dq\hat{f}(q)\hat{g}(-q),\qquad f,g\in L^2(\R)\cap L^1(\R),
\end{gather}
with the Fourier transform def\/ined by
\begin{gather}\label{Four}
\hat{h}(q)=\int_{\R}dp\exp(i\alpha pq)h(p),\qquad h=f,g.
\end{gather}
Rewriting $B$ as
\begin{gather}\label{Bnew}
B(b;x,y)=\int_{\R}dp\frac{G(p+(x-ib)/2)G(p-(x+ib)/2)\exp(i\alpha py)}{G(p-(x-ib)/2)G(p+(x+ib)/2)},\\
 b\in(0,2a),\qquad x,y\in\R,\nonumber
\end{gather}
we now def\/ine
\begin{gather}
f_1(p)\equiv\frac{G(p+(x-ib)/2)}{G(p-(x-ib)/2)},\qquad f_2(p)\equiv\frac{G(p+(x-ib)/2)}{G(p+(x+ib)/2)},
\\
g_1(p)\equiv\frac{G(p-(x+ib)/2)\exp(i\alpha py)}{G(p+(x+ib)/2)},\qquad
g_2(p)\equiv\frac{G(p-(x+ib)/2)\exp(i\alpha py)}{G(p-(x-ib)/2)}.
\end{gather}
We can calculate the Fourier transforms of these four functions by using Proposition~\ref{propositionC.1}. Doing so, we use the Plancherel relation~\eqref{Plan} and then replace $q$ by $z+y/2$ to obtain the two representations announced above:
\begin{gather}\label{rep2}
B(b;x,y)=G(\pm x +ia-ib)\int_{\R}dz \frac{G(z\pm (x-y)/2-ia+ib/2)}{G(z\pm (x+y)/2+ia-ib/2)},
\\
\label{rep3}
B(b;x,y)=G(ia-ib)^2\int_{\R}dz G(\pm z\pm y/2-ia+ib/2)\exp(i\alpha zx).
\end{gather}

Next, we compare~\eqref{rep3} to~\eqref{Bnew}, deducing that it can be rewritten as
\begin{gather}\label{BB}
B(b;x,y)=G(ia-ib)^2B(2a-b;y,x).
\end{gather}
Also, def\/ining a new function $C(b;x,y)$ by~\eqref{defC},
we see that~\eqref{rep2} amounts to
\begin{gather}\label{BC}
B(b;x,y)=\sqrt{\frac{2\pi}{\alpha}}G(\pm x +ia-ib)C(2a-b;x,y).
\end{gather}

The function $C(\aaa,b;x,y)$ is of central importance for what follows. Writing it as
\begin{gather}
C(\aaa,b;x,y)= \sqrt{\frac{\alpha}{2\pi}}\int_{\R}dz \prod_{\de=+,-}\frac{G(\aaa;z-u_{\de})}{G(\aaa;z-d_{\de})},
\end{gather}
where we have introduced
\begin{gather}
u_{\pm}\equiv \pm (x-y)/2+ib/2,\qquad d_{\pm}\equiv \pm (x+y)/2-ib/2,
\end{gather}
we infer that its behavior under analytic continuation in its 5 variables is immediate from the general analysis in Appendix~B of~I (with $N$ specialized to 2). We proceed to summarize the salient information.

To this end, we need the function $E(\aaa;z)$ discussed in~Appendix~\ref{appendixA}, cf.~\eqref{GE}. Introducing
\begin{gather}\label{P}
P(\aaa,b;x,y)\equiv C(\aaa,b;x,y)E(\aaa;\pm x+ib-ia)E(\aaa;\pm y+ib-ia),\!\!\!\!
\end{gather}
the product function $P(\aaa,b;x,y)$ extends from $(0,\infty)^2\times (0,a_{+}+a_-)\times \R^2$ to a function that is holomorphic in the domain
\begin{gather}
D(\aaa,b)\equiv \big\{ (\aaa,b,x,y)\in \C^5 \mid \re a_+>0,\ \re a_->0, \ \re (b/a_+a_-)>0\big\}.
\end{gather}
 Hence $C$ is meromorphic in $D(\aaa,b)$, with poles occurring solely at the zeros
\begin{gather}
\pm x=2ia-ib+z_{kl},\qquad \pm y=2ia-ib+z_{kl},\qquad k,l\in\N,
\end{gather}
of the $E$-product, cf.~\eqref{Eze}; moreover, the maximal multiplicity of a pole at $z=z_0$, with $z=x,y$, is given by the zero multiplicity at $z=z_0$ of the pertinent $E$-factor.

The corresponding meromorphy properties of $B$ are now clear from its relation~\eqref{BC} to $C$: It continues meromorphically to the domain $D(\aaa,2a-b)$. From~\eqref{BB} we then deduce that $B(\aaa,b;x,y)$ has a meromorphic extension to the larger domain $D_+$~\eqref{Dp}.
Using~\eqref{BC} again, it now follows that $C$ has a meromorphic extension to $D_+$ as well.

We proceed to obtain further information on the function $C$. First, we list features that are immediate from its def\/inition~\eqref{defC} and properties of the $G$-function, cf.~Appendix~\ref{appendixA}:
\begin{gather}\label{Cmodinv}
C(a_{+},a_{-},b;x,y)=C(a_{-},a_{+},b;x,y),\qquad  ({\rm modular\  invariance}),
\\
\label{Csd}
C(a_{+},a_{-},b;x,y)=C(a_{+},a_{-},b;y,x),\qquad  (\mbox{self-duality}),
\\
\label{Ceven}
C(b;x,y)=C(b;\de x, \de'y),\qquad \de,\de'=+,-,\qquad ({\rm evenness}),
\\
\label{Cscale}
C(\aaa,b;x,y)=C(\lambda  a_{+},\lambda a_{-}, \lambda b;\lambda x,\lambda y),\qquad \lambda>0,\qquad ({\rm scale\ invariance}),
\\
\label{Creal}
C(\aaa,b;x,y)\in\R,\!\!\!\!\qquad  \aaa>0,\!\!\!\!\qquad b\in(0,2a),\!\!\!\!\qquad x,y\in\R,\!\!\!\!\qquad \mbox{(real-valuedness)}.\!\!\!\!
\end{gather}
Clearly, the relations \eqref{Cmodinv}--\eqref{Cscale} are well def\/ined and hold true on $D_+$~\eqref{Dp}.  Also, the property~\eqref{Creal} can be rendered manifest by substituting the integral representation~\eqref{ghyp} in the four $G$-functions and combining factors to obtain
\begin{gather}
C(\aaa,b;x,y)=  \sqrt{\frac{\alpha}{2\pi}}  \int_{\R}dz\cos\left( \int_{\R}\frac{dw}{w}\frac{\sin(xw)\sin(yw)\cosh(bw)\sin(2zw)}{\sinh(a_{+}w)\sinh(a_-w)}\right)  \\
\phantom{C(\aaa,b;x,y)=}{}  \times \exp\left( \int_{\R}\frac{dw}{w}\left(\frac{\cos(xw)\cos(yw)\sinh(bw)\cos(2zw)}{\sinh(a_{+}w)\sinh(a_-w)}-\frac{b}{a_+a_-w}\right)\right),  \nonumber
\end{gather}
  where $(\aaa,b,x,y)\in(0,\infty)^2\times(0,2a)\times\R^2$.

Second, combining~\eqref{Csd} with~\eqref{BB} and~\eqref{BC}, we deduce
\begin{gather}\label{Cbsym}
\frac{G(ia-ib)C(b;x,y)}{G(x+ia-ib)G(y+ia-ib)}=
\frac{G(ib-ia)C(2a-b;x,y)}{G(x-ia+ib)G(y-ia+ib)},\qquad (b\mbox{-symmetry}).
\end{gather}
Third, we can use Proposition~\ref{propositionC.1} once more to obtain from~\eqref{defC} the explicit result
\begin{gather}\label{Cnorm}
C(b;x,ib)=G(ia-2ib)G(ib-ia)^2,\qquad ({\rm normalization}),
\end{gather}
where $(\aaa,b,x)\in(0,\infty)^2\times (0,a)\times \R$.
Last but not least, we claim that we have the joint eigenvalue equations
\begin{gather}\label{Cades}
A_{\pm}(b;x)C(b;x,y)=2c_{\pm}(y)C(b;x,y),\qquad A_{\pm}(b;y)C(b;x,y)=2c_{\pm}( x)C(b;x,y).
\end{gather}

 To prove this claim, we f\/irst note that the eigenvalue equations~\eqref{eigen} continue meromorphically to $D_+$. Next, we observe that the $G$-A$\De$Es~\eqref{Gades} imply
\begin{gather}
G(\pm x +ia-ib)^{-1}A_{\de}(b;x)G(\pm x +ia-ib)=A_{\de}(2a-b;x),\qquad \de=+,-.
\end{gather}
(Here, we view the l.h.s.\ as the product of three operators acting on meromorphic functions.)
Hence we obtain via~\eqref{BC}
\begin{gather}
A_{\pm}(2a-b;x)C(2a-b;x,y)=2c_{\pm}(y)C(2a-b;x,y),
\end{gather}
which is equivalent to the f\/irst two A$\De$Es in~\eqref{Cades}. The last two are then clear from the self-duality relation~\eqref{Csd}.

All of the properties of $C$ just derived also hold true for the function $G(ib-ia)\cR_r$ def\/ined by~\eqref{cRr} and~\eqref{cRdef}, cf.~Section~\ref{section2}. By using solely the eigenvalue properties~\eqref{Cades}, the evenness properties~\eqref{Ceven}, and the normalization~\eqref{Cnorm}, we can now show that these two functions coincide, as announced in the Introduction, cf.~\eqref{cRC}. Specif\/ically, applying the uniqueness argument in Subsection~\ref{section2.3} to $C$ in its dependence on $x$, we obtain the equality~\eqref{cRC} up to a~proportionality factor $p(\aaa,b,y)$. Repeating this argument for the $y$-dependence, we see that the proportionality factor can only depend on the parameters $a_+$, $a_-$, $b$. From the normalization relation~\eqref{Cnorm} it then follows that $p=1$, thus proving~\eqref{cRC}.

Let us now collect the resulting minimal representations of the $\cR$-function. From~\eqref{cRC} and~\eqref{cRr} we obtain
\begin{gather}\label{Ri}
\cR(b;x,y)=\frac{1}{\sqrt{a_+a_-}}\frac{G(2ib-ia)}{G(ib-ia)^2}
\int_{\R}dz \frac{G(z\pm (x-y)/2-ib/2)}{G(z\pm (x+y)/2+ib/2)}.
\end{gather}
Next, combining~\eqref{Bnew} and~\eqref{BC}, we deduce
\begin{gather}
\cR(b;x,y)=\frac{1}{\sqrt{a_+a_-}}\frac{G(2ib-ia)}{G(ib-ia)^2}
G(\pm x +ia-ib)\nonumber\\
\phantom{\cR(b;x,y)=}{} \times
\int_{\R}dz\frac{G(z\pm x/2-ia+ib/2)}{G(z\pm x/2+ia-ib/2)}\exp(i\alpha zy),\label{Rii}
\end{gather}
and using~\eqref{BB} we infer
\begin{gather}\label{Riii}
\cR(b;x,y)=\frac{1}{\sqrt{a_+a_-}}G(2ib-ia)
G(\pm x +ia-ib)
\int_{\R}dz\frac{G(z\pm y/2-ib/2)}{G(z\pm y/2+ib/2)}\exp(i\alpha zx).\!\!\!
\end{gather}
Finally, using the self-duality property of $\cR$, we obtain from~\eqref{Rii} and~\eqref{Riii} the representations
\begin{gather}
\cR(b;x,y)=\frac{1}{\sqrt{a_+a_-}}\frac{G(2ib-ia)}{G(ib-ia)^2}
G(\pm y +ia-ib)\nonumber\\
\phantom{\cR(b;x,y)=}{}\times
\int_{\R}dz\frac{G(z\pm y/2-ia+ib/2)}{G(z\pm y/2+ia-ib/2)}\exp(i\alpha zx),\label{Riv}
\\
\cR(b;x,y)=\frac{1}{\sqrt{a_+a_-}}G(2ib-ia)
G(\pm y +ia-ib)\nonumber\\
\phantom{\cR(b;x,y)=}{}\times
\int_{\R}dz\frac{G(z\pm x/2-ib/2)}{G(z\pm x/2+ib/2)}\exp(i\alpha zy).\label{Rv}
\end{gather}
The f\/ive representations~\eqref{Ri}--\eqref{Rv} are well def\/ined and hold true for $(\aaa,b,x,y)\in(0,\infty)^2\times(0,2a)\times\R^2$. (Combining~\eqref{Ri} with~\eqref{Cbsym}, we get a further representation that we do not consider.)

Taking stock of the above developments, we note that we might have started from the f\/irst minimal representation~\eqref{Ri} to {\em define} the $\cR$-function. Then many of its properties follow quite easily. On the other hand, it seems not feasible to give a direct proof of its crucial joint eigenfunction property. With hindsight, however, this can be shown by f\/irst obtaining the second representation~\eqref{Rii} (say) via Proposition~\ref{propositionC.1}, and then using the identity~\eqref{dv} to arrive at~\eqref{eigen}. From this the joint eigenfunction property~\eqref{Cades} follows as before.
